\section{Numerical algorithm and equilibrium audit}
\label{app:rewrite-algorithm}

\subsection{Numerical solution}

For the RSI-activation comparisons in
Sections~\ref{subsec:rewrite-rsi-half-decline} and~\ref{subsec:rewrite-rsi-research-share},
I compute numerical approximations to equilibrium trajectories by solving
a boundary-value problem (BVP). Initial capital $K_0$ and AI efficiency
$B_0$ are given, whereas initial consumption $C_0$ and the developer's
shadow price $q_0$ are determined jointly with the transition.
I solve the equilibrium differential equations over a finite horizon,
using the given stocks as initial conditions and the long-run
characterization in Proposition~\ref{prop:rewrite-equilibrium-regimes}
to construct terminal conditions.

For each $\sigma$, I first solve the existing-AI economy with efficiency
fixed at $B_0$ and research unavailable. A scalar root search in $\log K_0$
solves Equation~\eqref{eq:rewrite-rsi-initial-ratio} using the production
and monopoly equations. The resource constraint~\eqref{eq:rewrite-resource},
with $M=0$ and $\dot K=(n+\gamma)K$, gives pre-event consumption.
The post-activation problem inherits the stocks but selects consumption,
research, and the shadow price anew.

I use \texttt{scipy.integrate.solve\_bvp} from the open-source Python
library SciPy \citep{virtanenetal2020}. The solver implements fourth-order
collocation with residual control similar to that of
\citet{kierzenkashampine2001}. I use the computational coordinates
\[
 \left(\log K,\ \log\frac{B}{\overline B-B},\ \log C,\ \log q\right).
\]
The second coordinate expresses AI efficiency relative to the remaining
gap below its upper bound; it can grow without bound while AI efficiency
remains bounded. Subtracting the analytical long-run aggregate growth rate
times $t$ from each coordinate removes its asymptotic linear trend.
These transformations change the numerical representation, not the
economic variables or calendar time.
At each trial state the code solves the static monopoly problem and
recovers research from its first-order condition. The left boundary fixes
$K_0,B_0$; two right-boundary projections remove unstable components of
the normalized equilibrium dynamics near the long-run limit, using
Appendix~\ref{app:rewrite-finite-existence}.

For each of the eight parameter configurations, I start near its
analytical limiting point and move the initial stocks to their prescribed
values through a sequence of boundary-value problems. This continuation
changes the initial conditions of the numerical problem, not the dates
on the economic trajectory.

\phantomsection\label{app:rewrite-competitive-transition-numerics}
For the competitive-to-monopoly comparison in
Section~\ref{subsec:rewrite-competitive-transition}, I instead obtain
$K_0$ from $r=\rho+\gamma$ with marginal-cost pricing $p_X=1/B_0$;
at $\sigma=1.5$, $K_0=6.535$. Pre-event consumption follows from the same
resource constraint. I then solve the monopoly BVP at these initial
stocks, using the RSI-activation trajectories only as numerical initial
guesses, not as restrictions on consumption or research. This procedure
also applies to the four-elasticity comparisons in the
\href{https://oliverpardo1979.github.io/ai-growth-and-labor/\#app:rewrite-competitive-transition}{digital edition}.

\subsection{Accuracy and equilibrium checks}

I extend each computational horizon twice by 500 years and tighten the
collocation tolerance from $2\times10^{-6}$ to $10^{-8}$ and $10^{-9}$.
The final boundary tolerance is $10^{-11}$. The displayed window lies
inside each solved horizon; no spline is extrapolated.

The original equations are reconstructed independently with five-point
differences, with denser checks over the first ten years.
Acceptance of a computed trajectory requires dynamic residuals below $10^{-6}$, first-order
residuals below $10^{-9}$, final-horizon changes below $2\times10^{-5}$
in detrended logs, and terminal-coordinate gaps below $10^{-4}$.
The horizon comparison measures changes in the computed trajectory when
the terminal date is extended; the terminal-coordinate check measures its
distance from the analytical limiting point at that date.
I also check the pre-event equilibrium conditions and the continuity of
production quantities and prices at RSI activation.

The audit checks feasibility and developer global optimality. When
the concavity conditions of Lemma~\ref{lem:rewrite-hamiltonian-concavity}
fail, it checks the Hamiltonian
support condition in Lemma~\ref{lem:rewrite-developer-verification},
using counterfactual-state minimization and an analytical bound beyond
the grid. At each checked date, the support test compares the maximized
Hamiltonian at alternative AI-efficiency levels with its tangent at the
candidate state, holding the transformed shadow price and the dated
values of capital and effective labor fixed.
The analytical support bound below and the asymptotic rate $n-\rho<0$ of both transversality
expressions support the infinite-horizon continuation. These are
numerically verified equilibrium approximations, not interval-arithmetic
existence certificates for arbitrary initial stocks.
The competitive-to-monopoly paths satisfy the same residual,
horizon-refinement, transversality, and developer-optimality criteria.

\begin{table}[H]
\centering
\caption{Numerical accuracy of the illustrative RSI comparisons}
\label{tab:rewrite-rsi-illustrative-accuracy}
\small
\begin{tabular}{rrrrr}
\toprule
$\chi$ & $\sigma$ & Final horizon & Dynamic residual & Horizon change\\
\midrule
7.5 & 0.90 & 5,428.0 & $1.18\times10^{-9}$ & $3.06\times10^{-8}$\\
7.5 & 1.00 & 4,944.5 & $1.38\times10^{-9}$ & $3.46\times10^{-8}$\\
7.5 & 1.10 & 4,828.8 & $1.66\times10^{-9}$ & $3.73\times10^{-8}$\\
7.5 & 1.50 & 3,013.0 & $1.69\times10^{-9}$ & $1.56\times10^{-8}$\\
1.5 & 0.90 & 6,047.0 & $1.02\times10^{-9}$ & $4.30\times10^{-8}$\\
1.5 & 1.00 & 5,563.5 & $1.11\times10^{-9}$ & $3.80\times10^{-8}$\\
1.5 & 1.10 & 5,447.8 & $1.19\times10^{-9}$ & $4.08\times10^{-8}$\\
1.5 & 1.50 & 3,632.0 & $1.23\times10^{-9}$ & $1.80\times10^{-8}$\\
\bottomrule
\end{tabular}
\par\smallskip
\begin{minipage}{0.98\textwidth}\footnotesize
Horizons are measured in years. Dynamic residuals are the largest absolute
errors in the derivatives of the four logarithmic coordinates, across
the full-horizon and dense first-decade checks. Horizon change compares the
last two solutions in detrended logarithmic coordinates over years 0--500.
Reported error maxima are rounded upward. All eight paths pass the equation,
event-continuity, developer-optimality and long-run-continuation checks.
The $\sigma=1.50$ cases use the global Hamiltonian-support test and its
analytical long-run bound; the other cases pass the concavity test.
\end{minipage}
\end{table}


\subsection{Analytical checks on developer optimality}
\label{app:rewrite-numerical-support-checks}

At unit elasticity, define $\beta=(1-\alpha)\omega_X$. The static monopoly solution implies
$\pi_t(B)=a_tB^{\beta/(1-\beta)}$ for some $a_t>0$, conditional on the
dated values of capital and effective labor. Here $R_B$ is the transformed
AI-efficiency state in equation~\eqref{eq:rewrite-research-depth}.
Because $B(R_B)$ is
increasing and concave, $R_B\mapsto\pi_t(B(R_B))$ is globally concave whenever
$\beta\leq1/2$. The illustrative values $\beta=0.067$ and $\eta=0.20$
therefore satisfy both concavity conditions in
Lemma~\ref{lem:rewrite-hamiltonian-concavity}.

The AI-dominated continuation supplies an analytical long-run check on the
support inequality. For $1<\sigma\leq1/\alpha$, the service elasticity defined
in equation~\eqref{eq:rewrite-service-capability-elasticity} satisfies
$\varepsilon_{B,t}(b)\leq1/\alpha$ for every $t$ and $b$. Suppressing these
arguments, subtracting $\alpha(1-e_X)$ from the denominator in that equation
gives
\[
 (\sigma^{-1}-\alpha)(1-s_X)
 \left[(1-\sigma^{-1})(1-s_X)+(\sigma^{-1}-\alpha)s_X\right]\geq0.
\]
Choose $B_*\in(B_0,\overline B)$ with
$B_*/\overline B>\max\{0,(1-2\alpha)/(1-\alpha)\}$.
Operating profit is then concave in the transformed state $R_B$ for $B\geq B_*$ by
equation~\eqref{eq:rewrite-capped-profit-curvature}.
If $\eta\leq1/2$, the research term in
\eqref{eq:rewrite-maximized-depth-hamiltonian} is concave there as well.
For a candidate with $B_t\geq B_*$, it consequently suffices to check
\begin{equation}
 \widehat{\mathcal H}_t(R_{B,t})
 -\widehat{\mathcal H}'_t(R_{B,t})(R_{B,t}-R_{B,0})
 \geq\widehat{\mathcal H}_t(R_B(B_*)).
 \label{eq:rewrite-support-terminal-bound}
\end{equation}
Above $B_*$ the tangent bound follows from concavity. Below $B_*$ the
Hamiltonian is increasing in $B$, so its value is at most the right-hand
side, whereas the tangent is minimized at $R_{B,0}$.

Along an AI-dominated continuation, the difference between the two sides of
\eqref{eq:rewrite-support-terminal-bound}, divided by $Y_t$, converges to
\[
 \alpha(1-\alpha)
 \left[1-\left(\frac{B_*}{\overline B}\right)^{(1-\alpha)/\alpha}\right]>0.
\]
To obtain this limit, use $M/Y\to0$, $\psi(B)R_B\to0$, and
$\pi(B_*)/Y\to\alpha(1-\alpha)(B_*/\overline B)^{(1-\alpha)/\alpha}$.
Thus the global support condition holds sufficiently late on this
continuation even when the stronger curvature condition fails at low
counterfactual AI-efficiency levels. The finite transition still requires its own
support check.

\paragraph{Replication files.}

I use NumPy for array operations and linear algebra \citep{harrisetal2020}
and Matplotlib to generate the figures \citep{hunter2007}.
The code and simulation data are available in the
\href{https://github.com/oliverpardo1979/ai-growth-and-labor}{repository}.
The \href{https://github.com/oliverpardo1979/ai-growth-and-labor/blob/main/REPLICATION.md}
{replication guide} records package versions, implementation details,
numerical tolerances, and commands for checking and reproducing the
results, with links to the additional institutional experiments.
